\documentclass{chsmc}
\chsmcsetup{
	year = 2014,
	part = 1,
	type = B,
	show-answers = false,
}
\begin{document}

\section{填空题：本大题共 8 小题，每小题 8 分，共 64 分.}

\begin{enumerate}
	\item 若函数 $f(x)$ 的图象是由依次连接点 $(0,0)$, $(1,1)$ 和 $(2,3)$ 的折线，
		则 $f^{-1}(2) =$ \blankline
	\item 在正方体 $ABCD$-$A'B'C'D'$ 中，二面角 $A'$-$BD$-$C'$ 等于 \blankline
		（用反三角表示）
	\item 对于实数 $\mathbf{R}$ 的任意子集 $U$，我们在 $\mathbf{R}$ 上定义函数
		$f_U(x) = \begin{cases}
			1, & x \in U, \\
			0, & x \notin U.
		\end{cases}$
		如果 $A$, $B$ 是实数 $\mathbf{R}$ 的两个子集，则
		$f_A(x) + f_B(x) \equiv 1$ 的充分必要条件是 \blankline
	\item 若 $\triangle ABC$ 的三个内角 $A$, $B$, $C$ 的余切
		$\cot A$, $\cot B$, $\cot C$ 依次成等差数列，则角 $B$ 的最大值是 \blankline
	\item 实数列 $\{a_n\}$ 满足条件：$a_1 = \dfrac{\sqrt{2}}{2} + 1$，$a_2 = \sqrt{2} + 1$，
		$a_{n+1} + a_{n-1} = \dfrac{n}{a_n-a_{n-1}} + 2$（$n \geqslant 2$），
		则通项公式 $a_n =$ \blankline（$n \geqslant 1$）
	\item $F_1$, $F_2$ 是椭圆 $\dfrac{x^2}{a^2} + \dfrac{y^2}{b^2} = 1$（$a > b > 0$）的两个焦点，
		$P$ 为椭圆上的一点，如果 $\triangle PF_1F_2$ 的面积为 $1$，
		$\tan \angle PF_1F_2 = \dfrac{1}{2}$，$\tan \angle PF_2F_1 = -2$，
		则 $a =$ \blankline
	\item 将一副扑克牌中的大小王去掉，在剩下的 $52$ 张牌中随机地抽取 $5$ 张，
		其中至少有两张牌上的数字（或者字母 $J$, $Q$, $K$, $A$）相同的概率是 \blankline
		（要求计算出这个概率的数值，精确到 $0.01$）
	\item 设 $g(x) = \sqrt{x(1-x)}$ 是定义在区间 $[0,1]$ 上的函数，
		则函数 $y = xg(x)$ 的图象与 $x$ 轴所围成图形的面积是 \blankline
\end{enumerate}

\section{解答题：本大题共 3 个小题，满分 56 分. 解答应写出文字说明、证明过程或演算步骤.}

\begin{enumerate}[resume]
	\item (本题满分 16 分) 设数列 $\{a_n\}$ 的前 $n$ 项和 $S_n$ 组成的数列满足
		$S_n + S_{n+1} + S_{n+2} = 6n^2 + 9n + 7$（$n \geqslant 1$）.
		已知 $a_1 = 1$, $a_2 = 5$，求数列 $\{a_n\}$ 的通项公式.
	\item (本题满分 20 分) 设 $x_1$, $x_2$, $x_3$ 是多项式方程 $x^3 - 10x + 11 = 0$ 的三个根.
		\begin{enumerate}[(1)]
			\item 已知 $x_1$, $x_2$, $x_3$ 都落在区间 $(-5,5)$ 之中，求这三个根的整数部分;
			\item 证明：$\arctan x_1 + \arctan x_2 + \arctan x_3 = \dfrac{\pi}{4}$.
		\end{enumerate}
	\item (本题满分 20 分) 如下图，椭圆 $\Gamma\colon \dfrac{x^2}{4} + y^2 = 1$，
		$A(-2,0)$, $B(0,-1)$ 是椭圆 $\Gamma$ 上的两点，直线 $l_1\colon x = -2$, $l_2\colon y = -1$，
		$P(x_0,y_0)$（$x_0 > 0$, $y_0 > 0$）是 $\Gamma$ 上的一个动点，$l_3$ 是过点 $P$ 且与 $\Gamma$ 相切的直线，
		$C$、$D$、$E$ 分别是直线 $l_1$ 与 $l_2$, $l_2$ 与 $l_3$, $l_1$ 与 $l_3$ 的交点.
		求证：三条直线 $AD$, $BE$ 和 $CP$ 共点.
\end{enumerate}

\end{document}
